Take \(K=2\), \(d=1\), \(\Theta=[-1,1]\), constant contexts, and fixed potential outcomes \(Y_t(1)=Y_t(2)=1\). Define
\[
 g(1,x,y;\theta)=y-\theta/2,
 \qquad g(2,x,y;\theta)=-y-\theta/2.
\tag{1}
\]
Thus \(M_P(\theta)=-\theta\), so the unique causal moment root is \(\theta^*=0=E[Y_t(1)]-E[Y_t(2)]\) and \(H=-1\). The score and both derivatives are bounded on \(\Theta\), and the full-data array is deterministic, hence independent and identically distributed. The only randomness is the adaptive assignment.

At times \(t\le n\), assign arm one with probability \(1/2\). After observing \(S_n\), assign arm one at every time \(t>n\) with probability \(1/4\) if \(S_n>0\), and with probability \(1/2\) otherwise. This rule is \(\mathcal F_{t-1}\)-measurable, satisfies sequential randomization, and has overlap \(1/4\). At the root,
\[
 \xi_t=\frac{\mathbf1\{A_t=1\}}{p_t}
 -\frac{\mathbf1\{A_t=2\}}{1-p_t}.
\tag{2}
\]
Thus the conditional score variance is \(4\) when \(p_t=1/2\), and
\[
 \frac14(4)^2+\frac34(-4/3)^2=\frac{16}{3}
\tag{3}
\]
when \(p_t=1/4\). Consequently
\[
 \frac1T\langle S\rangle_T
 =\begin{cases}14/3,&S_n>0,\\4,&S_n\le0,
 \end{cases}
\tag{4}
\]
which is at least \(4\) on every assignment path.

To compute the exact limiting law, construct the first-block scores from independent fair signs and construct, independently of that block, two candidate second-block score arrays having respectively the \(p=1/4\) and \(p=1/2\) laws in (2). Taylor expansion of their joint characteristic functions at the origin gives
\[
 \left(\frac{S_n}{\sqrt n},
 \frac{S^{(1/4)}_{2n}-S_n}{\sqrt n},
 \frac{S^{(1/2)}_{2n}-S_n}{\sqrt n}
 \right)\Rightarrow
 \left(2Z,\frac4{\sqrt3}W_1,2W_2\right),
\tag{5}
\]
where \(Z,W_1,W_2\) are independent standard normals. Indeed, each one-step characteristic function equals \(1-u^2\sigma^2/(2n)+O(n^{-3/2})\) uniformly for fixed \(u\), and independence makes the joint characteristic function the product of the three limits. The boundary \(Z=0\) has probability zero, so selecting the second block by the sign of the first and applying the continuous mapping away from that null boundary gives
\[
 \frac{S_{2n}}{\sqrt n}\Rightarrow
 \mathbf1\{Z>0\}\left(2Z+\frac4{\sqrt3}W_1\right)
 +\mathbf1\{Z\le0\}(2Z+2W_2).
\tag{6}
\]

For realized QV, the variables \(\xi_t^2-E[\xi_t^2\mid\mathcal F_{t-1}]\) are bounded martingale differences. Orthogonality and Chebyshev's inequality therefore give
\[
 n^{-1}([S]_{2n}-\langle S\rangle_{2n})\longrightarrow0
\quad\text{in probability}.
\tag{7}
\]
Equations (4), (6), and (7) yield the displayed limit \(L\), with a single independent \(W\) representing the selected candidate block. Since the selected \(W\) has conditional mean zero and is independent of \(Z\),
\[
 \begin{aligned}
 E[L]
 &=\frac{2E[Z\mathbf1\{Z>0\}]}{\sqrt{28/3}}
 +\frac{2E[Z\mathbf1\{Z\le0\}]}{\sqrt8}\\
 &=\frac2{\sqrt{2\pi}}
 \left(\sqrt{\frac3{28}}-\frac1{\sqrt8}\right)\ne0.
 \end{aligned}
\tag{8}
\]
The two square roots differ, proving nonnormality as well as noncentering.

For completeness, the empirical moment is affine:
\[
 G_T(\theta)=T^{-1}\left(S_T-\frac\theta2R_T\right),
 \qquad
 R_T=\sum_{t=1}^T\left{
 \frac{\mathbf1\{A_t=1\}}{p_t}+
 \frac{\mathbf1\{A_t=2\}}{1-p_t}\right}.
\tag{9}
\]
The bounded martingale law gives \(R_T/T\to2\), because each summand has conditional mean \(2\). Hence the exact root \(\widehat\theta_T=2S_T/R_T\) satisfies \(\sqrt T\widehat\theta_T-S_T/\sqrt T\to0\) in probability. Its empirical derivative tends to \(-1\), and evaluating the bounded score at \(\widehat\theta_T\) changes the root-mean-square score by \(o_P(1)\). Choose \(c=1\) and \(\omega_{1,T}=4\) in the event-truncated sandwich. Equations (4) and (7), together with the preceding derivative and score comparisons, give
\[
 \widehat q_T(1)-T^{-1}[S]_T\longrightarrow0
 \quad\text{in probability}.
\tag{10}
\]
The oracle quantity converges to \(14/3\) on the positive regime and to \(4\) on the nonpositive regime, so the invertibility event and \(\widehat q_T(1)\ge\omega_{1,T}/2=2\) occur with probability tending to one. Therefore the feasible realized-QV IPW-Z pivot has the same limit \(L\). The later variance is measurable at time \(n\), but its selecting prefix is half the experiment, so the negligible-prefix condition of the repaired theorem fails exactly where this construction requires it.